
% Default to the notebook output style

    


% Inherit from the specified cell style.




    
\documentclass[11pt]{article}

    
    
    \usepackage[T1]{fontenc}
    % Nicer default font (+ math font) than Computer Modern for most use cases
    \usepackage{mathpazo}

    % Basic figure setup, for now with no caption control since it's done
    % automatically by Pandoc (which extracts ![](path) syntax from Markdown).
    \usepackage{graphicx}
    % We will generate all images so they have a width \maxwidth. This means
    % that they will get their normal width if they fit onto the page, but
    % are scaled down if they would overflow the margins.
    \makeatletter
    \def\maxwidth{\ifdim\Gin@nat@width>\linewidth\linewidth
    \else\Gin@nat@width\fi}
    \makeatother
    \let\Oldincludegraphics\includegraphics
    % Set max figure width to be 80% of text width, for now hardcoded.
    \renewcommand{\includegraphics}[1]{\Oldincludegraphics[width=.8\maxwidth]{#1}}
    % Ensure that by default, figures have no caption (until we provide a
    % proper Figure object with a Caption API and a way to capture that
    % in the conversion process - todo).
    \usepackage{caption}
    \DeclareCaptionLabelFormat{nolabel}{}
    \captionsetup{labelformat=nolabel}

    \usepackage{adjustbox} % Used to constrain images to a maximum size 
    \usepackage{xcolor} % Allow colors to be defined
    \usepackage{enumerate} % Needed for markdown enumerations to work
    \usepackage{geometry} % Used to adjust the document margins
    \usepackage{amsmath} % Equations
    \usepackage{amssymb} % Equations
    \usepackage{textcomp} % defines textquotesingle
    % Hack from http://tex.stackexchange.com/a/47451/13684:
    \AtBeginDocument{%
        \def\PYZsq{\textquotesingle}% Upright quotes in Pygmentized code
    }
    \usepackage{upquote} % Upright quotes for verbatim code
    \usepackage{eurosym} % defines \euro
    \usepackage[mathletters]{ucs} % Extended unicode (utf-8) support
    \usepackage[utf8x]{inputenc} % Allow utf-8 characters in the tex document
    \usepackage{fancyvrb} % verbatim replacement that allows latex
    \usepackage{grffile} % extends the file name processing of package graphics 
                         % to support a larger range 
    % The hyperref package gives us a pdf with properly built
    % internal navigation ('pdf bookmarks' for the table of contents,
    % internal cross-reference links, web links for URLs, etc.)
    \usepackage{hyperref}
    \usepackage{longtable} % longtable support required by pandoc >1.10
    \usepackage{booktabs}  % table support for pandoc > 1.12.2
    \usepackage[inline]{enumitem} % IRkernel/repr support (it uses the enumerate* environment)
    \usepackage[normalem]{ulem} % ulem is needed to support strikethroughs (\sout)
                                % normalem makes italics be italics, not underlines
    

    
    
    % Colors for the hyperref package
    \definecolor{urlcolor}{rgb}{0,.145,.698}
    \definecolor{linkcolor}{rgb}{.71,0.21,0.01}
    \definecolor{citecolor}{rgb}{.12,.54,.11}

    % ANSI colors
    \definecolor{ansi-black}{HTML}{3E424D}
    \definecolor{ansi-black-intense}{HTML}{282C36}
    \definecolor{ansi-red}{HTML}{E75C58}
    \definecolor{ansi-red-intense}{HTML}{B22B31}
    \definecolor{ansi-green}{HTML}{00A250}
    \definecolor{ansi-green-intense}{HTML}{007427}
    \definecolor{ansi-yellow}{HTML}{DDB62B}
    \definecolor{ansi-yellow-intense}{HTML}{B27D12}
    \definecolor{ansi-blue}{HTML}{208FFB}
    \definecolor{ansi-blue-intense}{HTML}{0065CA}
    \definecolor{ansi-magenta}{HTML}{D160C4}
    \definecolor{ansi-magenta-intense}{HTML}{A03196}
    \definecolor{ansi-cyan}{HTML}{60C6C8}
    \definecolor{ansi-cyan-intense}{HTML}{258F8F}
    \definecolor{ansi-white}{HTML}{C5C1B4}
    \definecolor{ansi-white-intense}{HTML}{A1A6B2}

    % commands and environments needed by pandoc snippets
    % extracted from the output of `pandoc -s`
    \providecommand{\tightlist}{%
      \setlength{\itemsep}{0pt}\setlength{\parskip}{0pt}}
    \DefineVerbatimEnvironment{Highlighting}{Verbatim}{commandchars=\\\{\}}
    % Add ',fontsize=\small' for more characters per line
    \newenvironment{Shaded}{}{}
    \newcommand{\KeywordTok}[1]{\textcolor[rgb]{0.00,0.44,0.13}{\textbf{{#1}}}}
    \newcommand{\DataTypeTok}[1]{\textcolor[rgb]{0.56,0.13,0.00}{{#1}}}
    \newcommand{\DecValTok}[1]{\textcolor[rgb]{0.25,0.63,0.44}{{#1}}}
    \newcommand{\BaseNTok}[1]{\textcolor[rgb]{0.25,0.63,0.44}{{#1}}}
    \newcommand{\FloatTok}[1]{\textcolor[rgb]{0.25,0.63,0.44}{{#1}}}
    \newcommand{\CharTok}[1]{\textcolor[rgb]{0.25,0.44,0.63}{{#1}}}
    \newcommand{\StringTok}[1]{\textcolor[rgb]{0.25,0.44,0.63}{{#1}}}
    \newcommand{\CommentTok}[1]{\textcolor[rgb]{0.38,0.63,0.69}{\textit{{#1}}}}
    \newcommand{\OtherTok}[1]{\textcolor[rgb]{0.00,0.44,0.13}{{#1}}}
    \newcommand{\AlertTok}[1]{\textcolor[rgb]{1.00,0.00,0.00}{\textbf{{#1}}}}
    \newcommand{\FunctionTok}[1]{\textcolor[rgb]{0.02,0.16,0.49}{{#1}}}
    \newcommand{\RegionMarkerTok}[1]{{#1}}
    \newcommand{\ErrorTok}[1]{\textcolor[rgb]{1.00,0.00,0.00}{\textbf{{#1}}}}
    \newcommand{\NormalTok}[1]{{#1}}
    
    % Additional commands for more recent versions of Pandoc
    \newcommand{\ConstantTok}[1]{\textcolor[rgb]{0.53,0.00,0.00}{{#1}}}
    \newcommand{\SpecialCharTok}[1]{\textcolor[rgb]{0.25,0.44,0.63}{{#1}}}
    \newcommand{\VerbatimStringTok}[1]{\textcolor[rgb]{0.25,0.44,0.63}{{#1}}}
    \newcommand{\SpecialStringTok}[1]{\textcolor[rgb]{0.73,0.40,0.53}{{#1}}}
    \newcommand{\ImportTok}[1]{{#1}}
    \newcommand{\DocumentationTok}[1]{\textcolor[rgb]{0.73,0.13,0.13}{\textit{{#1}}}}
    \newcommand{\AnnotationTok}[1]{\textcolor[rgb]{0.38,0.63,0.69}{\textbf{\textit{{#1}}}}}
    \newcommand{\CommentVarTok}[1]{\textcolor[rgb]{0.38,0.63,0.69}{\textbf{\textit{{#1}}}}}
    \newcommand{\VariableTok}[1]{\textcolor[rgb]{0.10,0.09,0.49}{{#1}}}
    \newcommand{\ControlFlowTok}[1]{\textcolor[rgb]{0.00,0.44,0.13}{\textbf{{#1}}}}
    \newcommand{\OperatorTok}[1]{\textcolor[rgb]{0.40,0.40,0.40}{{#1}}}
    \newcommand{\BuiltInTok}[1]{{#1}}
    \newcommand{\ExtensionTok}[1]{{#1}}
    \newcommand{\PreprocessorTok}[1]{\textcolor[rgb]{0.74,0.48,0.00}{{#1}}}
    \newcommand{\AttributeTok}[1]{\textcolor[rgb]{0.49,0.56,0.16}{{#1}}}
    \newcommand{\InformationTok}[1]{\textcolor[rgb]{0.38,0.63,0.69}{\textbf{\textit{{#1}}}}}
    \newcommand{\WarningTok}[1]{\textcolor[rgb]{0.38,0.63,0.69}{\textbf{\textit{{#1}}}}}
    
    
    % Define a nice break command that doesn't care if a line doesn't already
    % exist.
    \def\br{\hspace*{\fill} \\* }
    % Math Jax compatability definitions
    \def\gt{>}
    \def\lt{<}
    % Document parameters
    \title{P2\_2019\_1}
    
    
    

    % Pygments definitions
    
\makeatletter
\def\PY@reset{\let\PY@it=\relax \let\PY@bf=\relax%
    \let\PY@ul=\relax \let\PY@tc=\relax%
    \let\PY@bc=\relax \let\PY@ff=\relax}
\def\PY@tok#1{\csname PY@tok@#1\endcsname}
\def\PY@toks#1+{\ifx\relax#1\empty\else%
    \PY@tok{#1}\expandafter\PY@toks\fi}
\def\PY@do#1{\PY@bc{\PY@tc{\PY@ul{%
    \PY@it{\PY@bf{\PY@ff{#1}}}}}}}
\def\PY#1#2{\PY@reset\PY@toks#1+\relax+\PY@do{#2}}

\expandafter\def\csname PY@tok@w\endcsname{\def\PY@tc##1{\textcolor[rgb]{0.73,0.73,0.73}{##1}}}
\expandafter\def\csname PY@tok@c\endcsname{\let\PY@it=\textit\def\PY@tc##1{\textcolor[rgb]{0.25,0.50,0.50}{##1}}}
\expandafter\def\csname PY@tok@cp\endcsname{\def\PY@tc##1{\textcolor[rgb]{0.74,0.48,0.00}{##1}}}
\expandafter\def\csname PY@tok@k\endcsname{\let\PY@bf=\textbf\def\PY@tc##1{\textcolor[rgb]{0.00,0.50,0.00}{##1}}}
\expandafter\def\csname PY@tok@kp\endcsname{\def\PY@tc##1{\textcolor[rgb]{0.00,0.50,0.00}{##1}}}
\expandafter\def\csname PY@tok@kt\endcsname{\def\PY@tc##1{\textcolor[rgb]{0.69,0.00,0.25}{##1}}}
\expandafter\def\csname PY@tok@o\endcsname{\def\PY@tc##1{\textcolor[rgb]{0.40,0.40,0.40}{##1}}}
\expandafter\def\csname PY@tok@ow\endcsname{\let\PY@bf=\textbf\def\PY@tc##1{\textcolor[rgb]{0.67,0.13,1.00}{##1}}}
\expandafter\def\csname PY@tok@nb\endcsname{\def\PY@tc##1{\textcolor[rgb]{0.00,0.50,0.00}{##1}}}
\expandafter\def\csname PY@tok@nf\endcsname{\def\PY@tc##1{\textcolor[rgb]{0.00,0.00,1.00}{##1}}}
\expandafter\def\csname PY@tok@nc\endcsname{\let\PY@bf=\textbf\def\PY@tc##1{\textcolor[rgb]{0.00,0.00,1.00}{##1}}}
\expandafter\def\csname PY@tok@nn\endcsname{\let\PY@bf=\textbf\def\PY@tc##1{\textcolor[rgb]{0.00,0.00,1.00}{##1}}}
\expandafter\def\csname PY@tok@ne\endcsname{\let\PY@bf=\textbf\def\PY@tc##1{\textcolor[rgb]{0.82,0.25,0.23}{##1}}}
\expandafter\def\csname PY@tok@nv\endcsname{\def\PY@tc##1{\textcolor[rgb]{0.10,0.09,0.49}{##1}}}
\expandafter\def\csname PY@tok@no\endcsname{\def\PY@tc##1{\textcolor[rgb]{0.53,0.00,0.00}{##1}}}
\expandafter\def\csname PY@tok@nl\endcsname{\def\PY@tc##1{\textcolor[rgb]{0.63,0.63,0.00}{##1}}}
\expandafter\def\csname PY@tok@ni\endcsname{\let\PY@bf=\textbf\def\PY@tc##1{\textcolor[rgb]{0.60,0.60,0.60}{##1}}}
\expandafter\def\csname PY@tok@na\endcsname{\def\PY@tc##1{\textcolor[rgb]{0.49,0.56,0.16}{##1}}}
\expandafter\def\csname PY@tok@nt\endcsname{\let\PY@bf=\textbf\def\PY@tc##1{\textcolor[rgb]{0.00,0.50,0.00}{##1}}}
\expandafter\def\csname PY@tok@nd\endcsname{\def\PY@tc##1{\textcolor[rgb]{0.67,0.13,1.00}{##1}}}
\expandafter\def\csname PY@tok@s\endcsname{\def\PY@tc##1{\textcolor[rgb]{0.73,0.13,0.13}{##1}}}
\expandafter\def\csname PY@tok@sd\endcsname{\let\PY@it=\textit\def\PY@tc##1{\textcolor[rgb]{0.73,0.13,0.13}{##1}}}
\expandafter\def\csname PY@tok@si\endcsname{\let\PY@bf=\textbf\def\PY@tc##1{\textcolor[rgb]{0.73,0.40,0.53}{##1}}}
\expandafter\def\csname PY@tok@se\endcsname{\let\PY@bf=\textbf\def\PY@tc##1{\textcolor[rgb]{0.73,0.40,0.13}{##1}}}
\expandafter\def\csname PY@tok@sr\endcsname{\def\PY@tc##1{\textcolor[rgb]{0.73,0.40,0.53}{##1}}}
\expandafter\def\csname PY@tok@ss\endcsname{\def\PY@tc##1{\textcolor[rgb]{0.10,0.09,0.49}{##1}}}
\expandafter\def\csname PY@tok@sx\endcsname{\def\PY@tc##1{\textcolor[rgb]{0.00,0.50,0.00}{##1}}}
\expandafter\def\csname PY@tok@m\endcsname{\def\PY@tc##1{\textcolor[rgb]{0.40,0.40,0.40}{##1}}}
\expandafter\def\csname PY@tok@gh\endcsname{\let\PY@bf=\textbf\def\PY@tc##1{\textcolor[rgb]{0.00,0.00,0.50}{##1}}}
\expandafter\def\csname PY@tok@gu\endcsname{\let\PY@bf=\textbf\def\PY@tc##1{\textcolor[rgb]{0.50,0.00,0.50}{##1}}}
\expandafter\def\csname PY@tok@gd\endcsname{\def\PY@tc##1{\textcolor[rgb]{0.63,0.00,0.00}{##1}}}
\expandafter\def\csname PY@tok@gi\endcsname{\def\PY@tc##1{\textcolor[rgb]{0.00,0.63,0.00}{##1}}}
\expandafter\def\csname PY@tok@gr\endcsname{\def\PY@tc##1{\textcolor[rgb]{1.00,0.00,0.00}{##1}}}
\expandafter\def\csname PY@tok@ge\endcsname{\let\PY@it=\textit}
\expandafter\def\csname PY@tok@gs\endcsname{\let\PY@bf=\textbf}
\expandafter\def\csname PY@tok@gp\endcsname{\let\PY@bf=\textbf\def\PY@tc##1{\textcolor[rgb]{0.00,0.00,0.50}{##1}}}
\expandafter\def\csname PY@tok@go\endcsname{\def\PY@tc##1{\textcolor[rgb]{0.53,0.53,0.53}{##1}}}
\expandafter\def\csname PY@tok@gt\endcsname{\def\PY@tc##1{\textcolor[rgb]{0.00,0.27,0.87}{##1}}}
\expandafter\def\csname PY@tok@err\endcsname{\def\PY@bc##1{\setlength{\fboxsep}{0pt}\fcolorbox[rgb]{1.00,0.00,0.00}{1,1,1}{\strut ##1}}}
\expandafter\def\csname PY@tok@kc\endcsname{\let\PY@bf=\textbf\def\PY@tc##1{\textcolor[rgb]{0.00,0.50,0.00}{##1}}}
\expandafter\def\csname PY@tok@kd\endcsname{\let\PY@bf=\textbf\def\PY@tc##1{\textcolor[rgb]{0.00,0.50,0.00}{##1}}}
\expandafter\def\csname PY@tok@kn\endcsname{\let\PY@bf=\textbf\def\PY@tc##1{\textcolor[rgb]{0.00,0.50,0.00}{##1}}}
\expandafter\def\csname PY@tok@kr\endcsname{\let\PY@bf=\textbf\def\PY@tc##1{\textcolor[rgb]{0.00,0.50,0.00}{##1}}}
\expandafter\def\csname PY@tok@bp\endcsname{\def\PY@tc##1{\textcolor[rgb]{0.00,0.50,0.00}{##1}}}
\expandafter\def\csname PY@tok@fm\endcsname{\def\PY@tc##1{\textcolor[rgb]{0.00,0.00,1.00}{##1}}}
\expandafter\def\csname PY@tok@vc\endcsname{\def\PY@tc##1{\textcolor[rgb]{0.10,0.09,0.49}{##1}}}
\expandafter\def\csname PY@tok@vg\endcsname{\def\PY@tc##1{\textcolor[rgb]{0.10,0.09,0.49}{##1}}}
\expandafter\def\csname PY@tok@vi\endcsname{\def\PY@tc##1{\textcolor[rgb]{0.10,0.09,0.49}{##1}}}
\expandafter\def\csname PY@tok@vm\endcsname{\def\PY@tc##1{\textcolor[rgb]{0.10,0.09,0.49}{##1}}}
\expandafter\def\csname PY@tok@sa\endcsname{\def\PY@tc##1{\textcolor[rgb]{0.73,0.13,0.13}{##1}}}
\expandafter\def\csname PY@tok@sb\endcsname{\def\PY@tc##1{\textcolor[rgb]{0.73,0.13,0.13}{##1}}}
\expandafter\def\csname PY@tok@sc\endcsname{\def\PY@tc##1{\textcolor[rgb]{0.73,0.13,0.13}{##1}}}
\expandafter\def\csname PY@tok@dl\endcsname{\def\PY@tc##1{\textcolor[rgb]{0.73,0.13,0.13}{##1}}}
\expandafter\def\csname PY@tok@s2\endcsname{\def\PY@tc##1{\textcolor[rgb]{0.73,0.13,0.13}{##1}}}
\expandafter\def\csname PY@tok@sh\endcsname{\def\PY@tc##1{\textcolor[rgb]{0.73,0.13,0.13}{##1}}}
\expandafter\def\csname PY@tok@s1\endcsname{\def\PY@tc##1{\textcolor[rgb]{0.73,0.13,0.13}{##1}}}
\expandafter\def\csname PY@tok@mb\endcsname{\def\PY@tc##1{\textcolor[rgb]{0.40,0.40,0.40}{##1}}}
\expandafter\def\csname PY@tok@mf\endcsname{\def\PY@tc##1{\textcolor[rgb]{0.40,0.40,0.40}{##1}}}
\expandafter\def\csname PY@tok@mh\endcsname{\def\PY@tc##1{\textcolor[rgb]{0.40,0.40,0.40}{##1}}}
\expandafter\def\csname PY@tok@mi\endcsname{\def\PY@tc##1{\textcolor[rgb]{0.40,0.40,0.40}{##1}}}
\expandafter\def\csname PY@tok@il\endcsname{\def\PY@tc##1{\textcolor[rgb]{0.40,0.40,0.40}{##1}}}
\expandafter\def\csname PY@tok@mo\endcsname{\def\PY@tc##1{\textcolor[rgb]{0.40,0.40,0.40}{##1}}}
\expandafter\def\csname PY@tok@ch\endcsname{\let\PY@it=\textit\def\PY@tc##1{\textcolor[rgb]{0.25,0.50,0.50}{##1}}}
\expandafter\def\csname PY@tok@cm\endcsname{\let\PY@it=\textit\def\PY@tc##1{\textcolor[rgb]{0.25,0.50,0.50}{##1}}}
\expandafter\def\csname PY@tok@cpf\endcsname{\let\PY@it=\textit\def\PY@tc##1{\textcolor[rgb]{0.25,0.50,0.50}{##1}}}
\expandafter\def\csname PY@tok@c1\endcsname{\let\PY@it=\textit\def\PY@tc##1{\textcolor[rgb]{0.25,0.50,0.50}{##1}}}
\expandafter\def\csname PY@tok@cs\endcsname{\let\PY@it=\textit\def\PY@tc##1{\textcolor[rgb]{0.25,0.50,0.50}{##1}}}

\def\PYZbs{\char`\\}
\def\PYZus{\char`\_}
\def\PYZob{\char`\{}
\def\PYZcb{\char`\}}
\def\PYZca{\char`\^}
\def\PYZam{\char`\&}
\def\PYZlt{\char`\<}
\def\PYZgt{\char`\>}
\def\PYZsh{\char`\#}
\def\PYZpc{\char`\%}
\def\PYZdl{\char`\$}
\def\PYZhy{\char`\-}
\def\PYZsq{\char`\'}
\def\PYZdq{\char`\"}
\def\PYZti{\char`\~}
% for compatibility with earlier versions
\def\PYZat{@}
\def\PYZlb{[}
\def\PYZrb{]}
\makeatother


    % Exact colors from NB
    \definecolor{incolor}{rgb}{0.0, 0.0, 0.5}
    \definecolor{outcolor}{rgb}{0.545, 0.0, 0.0}



    
    % Prevent overflowing lines due to hard-to-break entities
    \sloppy 
    % Setup hyperref package
    \hypersetup{
      breaklinks=true,  % so long urls are correctly broken across lines
      colorlinks=true,
      urlcolor=urlcolor,
      linkcolor=linkcolor,
      citecolor=citecolor,
      }
    % Slightly bigger margins than the latex defaults
    
    \geometry{verbose,tmargin=1in,bmargin=1in,lmargin=1in,rmargin=1in}
    
    

    \begin{document}
    
    
    \maketitle
    
    

    
    \subsubsection{\texorpdfstring{Universidade Federal de Santa Maria,
Departamento de Física FSC 1046 - Mecânica Quântica II, Prof. Jonas
Maziero Semestre 2019/1, Data: 12 de julho de
2019}{Universidade Federal de Santa Maria, Departamento de Física   FSC 1046 - Mecânica Quântica II, Prof. Jonas Maziero   Semestre 2019/1, Data: 12 de julho de 2019}}\label{universidade-federal-de-santa-maria-departamento-de-fuxedsica-fsc-1046---mecuxe2nica-quuxe2ntica-ii-prof.-jonas-maziero-semestre-20191-data-12-de-julho-de-2019}

\section{Prova II}\label{prova-ii}

    \(1.\) Considere a decomposição espectral do Hamiltoniano da interação
dipolar magnética:

\begin{equation}
H_{d} = 0|\Psi_{-}\rangle\langle\Psi_{-}|+2D|\Psi_{+}\rangle\langle\Psi_{+}|-D(|\Phi_{-}\rangle\langle\Phi_{-}|+|\Phi_{+}
\rangle\langle\Phi_{+}|),
\end{equation}

em que
\(|\Psi_{\pm}\rangle=2^{-1/2}(|z_{+}\rangle\otimes|z_{-}\rangle\pm|z_{-}\rangle\otimes|z_{+}\rangle)\)
e
\(|\Phi_{\pm}\rangle=2^{-1/2}(|z_{+}\rangle\otimes|z_{+}\rangle\pm|z_{-}\rangle\otimes|z_{-}\rangle)\).
Para duas partículas de spin 1/2 preparadas no estado
\(|\psi_{0}\rangle = |z_{+}\rangle\otimes|z_{-}\rangle\), interagindo
segundo \(H_{d}\), verifique que o estado evoluído desse sistema em
\(t=\pi\hbar/4D\) pode ser escrito como

\begin{equation}
e^{-i\pi/4}\left(|z_{+}\rangle\otimes|z_{-}\rangle - i|z_{-}\rangle\otimes|z_{+}\rangle\right)/\sqrt{2}.
\end{equation}

    \emph{Sol.} Como o Hamiltoniano é independente do tempo, o operador de
evolução temporal pode ser esrito como:

\begin{equation}
U_{t} = e^{-iHt/\hbar} = |\Psi_{-}\rangle\langle\Psi_{-}|+e^{-i2Dt/\hbar}|\Psi_{+}\rangle\langle\Psi_{+}|+e^{iDt/\hbar}(|\Phi_{-}\rangle\langle\Phi_{-}|+|\Phi_{+}\rangle\langle\Phi_{+}|).
\end{equation}

Usando que \(\langle\psi_{0}|\Psi_{\pm}\rangle=1/\sqrt{2}\) e que
\(\langle\psi_{0}|\Phi_{\pm}\rangle=0\) teremos

\begin{align}
|\psi_{t}\rangle & = U_{t}|\psi_{0}\rangle \\
& = |\Psi_{-}\rangle\langle\Psi_{-}|\psi_{0}\rangle+e^{-i2Dt/\hbar}|\Psi_{+}\rangle\langle\Psi_{+}|\psi_{0}\rangle+e^{iDt/\hbar}(|\Phi_{-}\rangle\langle\Phi_{-}|\psi_{0}\rangle+|\Phi_{+}\rangle\langle\Phi_{+}|\psi_{0}\rangle) \\
&  = e^{-iDt/\hbar}\left(e^{iDt/\hbar}|\Psi_{-}\rangle + e^{-iDt/\hbar}|\Psi_{+}\rangle\right)/\sqrt{2} \\
&  = e^{-iDt/\hbar}\left((e^{iDt/\hbar}+e^{-iDt/\hbar})|z_{+}\rangle\otimes|z_{-}\rangle + (-e^{iDt/\hbar}+e^{-iDt/\hbar})|z_{-}\rangle\otimes|z_{+}\rangle \right)/2 \\
&  = e^{-iDt/\hbar}\left(\cos(Dt/\hbar)|z_{+}\rangle\otimes|z_{-}\rangle - i\sin(Dt/\hbar)|z_{-}\rangle\otimes|z_{+}\rangle \right).
\end{align}

Fazendo \(Dt/\hbar=\pi/4\) obteremos o resultado esperado.

    \(2.\) Partindo da equação de Schrödinger
\(i\hbar\partial_{t}|\psi_{t}\rangle=H|\psi_{t}\rangle\) com
\(H=P^{2}/2m + V_{t}\), assumindo que
\(V_{t}|\vec{x}\rangle=u_{t}(\vec{x})|\vec{x}\rangle\) com
\(u_{t}(\vec{x})\) sendo uma função escalar, e dado que
\(\langle x_{j}|P_{j}|\psi\rangle = \frac{\hbar}{i}\partial_{x_{j}}\langle x_{j}|\psi\rangle\),
obtenha a equação de Schödinger na base de posição:

\begin{equation}
i\hbar\partial_{t}\psi_{t}(\vec{x}) = \frac{-\hbar^{2}}{2m}\nabla^{2}_{\vec{x}}\psi_{t}(\vec{x}) + u_{t}(\vec{x})\psi_{t}(\vec{x}).
\end{equation}

    \emph{Sol.} Está nas notas de aula: " Seguindo, tomemos o produto
interno de \(|\vec{x}\rangle\) com a eq. de Schrödinger:

\begin{equation}
i\hbar\partial_{t}\langle\vec{x}|\psi_{t}\rangle = \langle\vec{x}|\frac{P^{2}}{2m}|\psi_{t}\rangle + \langle\vec{x}|U_{t}|\psi_{t}\rangle.
\end{equation}

Agora,

\begin{align}
\langle x_{j}|P_{j}^{2}|\psi\rangle & = \int dx_{j}'\langle x_{j}|P_{j}|x_{j}'\rangle\langle x_{j}'|P_{j}|\psi\rangle = \int dx_{j}'\langle x_{j}|P_{j}|x_{j}'\rangle(\hbar/i)(\partial_{x_{j}'}\psi(x_{j}'))  \\
& = \int dx_{j}'\langle x_{j}|P_{j}\left((\hbar/i)(\partial_{x_{j}'}\psi(x_{j}'))|x_{j}'\rangle\right) = \int dx_{j}'(\hbar/i)\partial_{x_{j}}\left(\langle x_{j}|(\hbar/i)(\partial_{x_{j}'}\psi(x_{j}'))|x_{j}'\rangle\right) \\ 
& = -\hbar^{2}\int dx_{j}'\partial_{x_{j}}\left((\partial_{x_{j}'}\psi(x_{j}'))\langle x_{j}|x_{j}'\rangle\right) = -\hbar^{2}\partial_{x_{j}}\int dx_{j}'(\partial_{x_{j}'}\psi(x_{j}'))\delta(x_{j}-x_{j}') \\
& = -\hbar^{2}\partial_{x_{j}}(\partial_{x_{j}}\psi(x_{j})) = -\hbar^{2}\partial_{x_{j}x_{j}}\psi(x_{j}).
\end{align}

Ainda, usando
\(\langle\vec{x}|P_{j}^{2}|\psi\rangle=-\hbar^{2}\partial_{x_{j}x_{j}}\psi(\vec{x})\),
vem que

\begin{align}
\langle\vec{x}|P^{2}|\psi\rangle & = \langle\vec{x}|\sum_{j}P_{j}^{2}|\psi\rangle = \sum_{j}\langle\vec{x}|P_{j}^{2}|\psi\rangle = -\hbar^{2}\sum_{j}\partial_{x_{j}x_{j}}\psi(\vec{x}) \\
     & = -\hbar^{2}\nabla^{2}_{\vec{x}}\psi(\vec{x}).
\end{align}

Substituindo na eq. de Schrödinger acima, e usando

\begin{equation}
U_{t}|\vec{x}\rangle = u_{t}(\vec{x})|\vec{x}\rangle
\end{equation}

com \(u_{t}(\vec{x})\) sendo uma função escalar, obtemos a equação de
Schrödinger para os coeficientes
\(\left(c_{\vec{x}}=\langle\vec{x}|\psi_{t}\rangle=\psi_{t}(\vec{x})\right)\)
do estado na base de posição
\(\left(|\psi_{t}\rangle=\int d^{3}x c_{\vec{x}}|\vec{x}\rangle\right)\):

\begin{equation}
i\hbar\partial_{t}\psi_{t}(\vec{x}) = \frac{-\hbar^{2}}{2m}\nabla^{2}_{\vec{x}}\psi_{t}(\vec{x}) + u_{t}(\vec{x})\psi_{t}(\vec{x}).
\end{equation}

"

    \(3.\) Considere o problema da soma de momento angular em Mecânica
Quântica. Obtenha a base \(|j_{a},j_{b},j,m\rangle\) em termos da base
\(|j_{a},m_{a}\rangle\otimes|j_{b},m_{b}\rangle\) para \(j_{a}=1/2\) e
\(j_{b}=1/2\).

    \emph{Sol.} Temos que \(m_{a}=-1/2,1/2\), \(m_{b}=-1/2,1/2\),
\(j=0\Rightarrow m=0\) e \(j=1\Rightarrow m=-1,0,1\). Temos a restrição
\(m\ne m_{a}+m_{b}\Rightarrow C^{j_{a}j_{b}j}_{m_{a}m_{b}m}=0\). Assim

\begin{equation}
|j=1,m=1\rangle = |m_{a}=1/2\rangle\otimes|m_{b}=1/2\rangle.
\end{equation}

Aplicando o operador de abaixamento,

\begin{equation}
J_{-}|j=1,m=1\rangle = \hbar\sqrt{1(1+1)-1(1-1)}|j=1,m=0\rangle = \hbar\sqrt{2}|j=1,m=0\rangle.
\end{equation}

Assim,

\begin{align}
|j=1,m=0\rangle & = \hbar^{-1}2^{-1/2}\left( J_{-}^{a}|m_{a}=1/2\rangle\otimes|m_{b}=1/2\rangle + |m_{a}=1/2\rangle\otimes J_{-}^{a}|m_{b}=1/2\rangle \right) \\
 & = \hbar^{-1}2^{-1/2}\left( \hbar\sqrt{1/2(1/2+1)-1/2(1/2-1)}|m_{a}=-1/2\rangle\otimes|m_{b}=1/2\rangle + |m_{a}=1/2\rangle\otimes \hbar\sqrt{1/2(1/2+1)-1/2(1/2-1)}|m_{b}=-1/2\rangle \right) \\
 & = 2^{-1/2}\left(|m_{a}=-1/2\rangle\otimes|m_{b}=1/2\rangle + |m_{a}=1/2\rangle\otimes |m_{b}=-1/2\rangle \right).
\end{align}

Como devemos ter \(m=m_{a}+m_{b}\), segue que

\begin{equation}
|j=1,m=-1\rangle = |m_{a}=-1/2\rangle\otimes|m_{b}=-1/2\rangle.
\end{equation}

Para o outro valor de \(j\) (\(j=0\)), precisaremos usar a normalização
do estado

\begin{equation}
|j=0,m=0\rangle = C^{1/2,1/2,0}_{-1/2,1/2,0}|m_{a}=-1/2\rangle\otimes|m_{b}=1/2\rangle + C^{1/2,1/2,0}_{1/2,-1/2,0}|m_{a}=1/2\rangle\otimes |m_{b}=-1/2\rangle,
\end{equation}

que nos fornece a equação:

\begin{equation}
|C^{1/2,1/2,0}_{-1/2,1/2,0}|^{2}+|C^{1/2,1/2,0}_{1/2,-1/2,0}|^{2}=1.
\end{equation}

Também precisamos utilizar a ortogonalidade com outros estados com
diferentes valores de \(j\) e mesmo valor de \(m\):

\begin{align}
\langle j=1,m=0|j=0,m=0\rangle = 0 = 2^{-1/2}C^{1/2,1/2,0}_{-1/2,1/2,0} + 2^{-1/2}C^{1/2,1/2,0}_{1/2,-1/2,0} \Rightarrow C^{1/2,1/2,0}_{-1/2,1/2,0} = -C^{1/2,1/2,0}_{1/2,-1/2,0}.
\end{align}

Assumindo coeficientes reais teremos que
\(C^{1/2,1/2,0}_{-1/2,1/2,0}=-C^{1/2,1/2,0}_{1/2,-1/2,0}=2^{-1/2}\) e

\begin{equation}
|j=0,m=0\rangle = 2^{-1/2}\left(|m_{a}=-1/2\rangle\otimes|m_{b}=1/2\rangle - |m_{a}=1/2\rangle\otimes |m_{b}=-1/2\rangle \right).
\end{equation}

    \(4.\) Considere o Hamiltoniano \(H=H_{0}+H_{1}\) com
\(H_{0}=-\omega\sigma_{y}\) e Hamiltoniano perturbativo
\(H_{1}=-\epsilon\sigma_{x}\), i.e., \(\omega\gg\epsilon\). Calcule as
correções perturbativas de primeira e de segunda ordem para as
auto-energias desse sistema e as correções de primeira ordem para os
autovetores de \(H\).

    \emph{Sol.} Temos
\(H_{0}=-\omega|y_{+}\rangle\langle y_{+}|+\omega|y_{-}\rangle\langle y_{-}|\).
Assim \(E_{1}^{(0)}=-\omega\), \(E_{2}^{(0)}=\omega\),
\(|\psi_{1}^{(0)}\rangle = |y_{+}\rangle\) e
\(|\psi_{2}^{(0)}\rangle = |y_{-}\rangle\). Utilizaremos a relação

\begin{equation}
\sigma_{x}|y_{\pm}\rangle = \sigma_{x}(|z_{+}\rangle \pm i|z_{-}\rangle)/\sqrt{2} = (\sigma_{x}|z_{+}\rangle \pm i\sigma_{x}|z_{-}\rangle)/\sqrt{2} = (|z_{-}\rangle \pm i|z_{+}\rangle)/\sqrt{2} = \pm i|y_{\mp}\rangle.
\end{equation}

Correções de primeira ordem para a energia
(\(E_{n}^{(1)} = \langle\psi_{n}^{(0)}|H_{1}|\psi_{n}^{(0)}\rangle\)):

\begin{align}
E_{1}^{(1)} & = \langle\psi_{1}^{(0)}|H_{1}|\psi_{1}^{(0)}\rangle = -\epsilon\langle y_{+}|\sigma_{x}|y_{+}\rangle = -i\epsilon\langle y_{+}|y_{-}\rangle = 0, \\
E_{2}^{(1)} & = \langle\psi_{2}^{(0)}|H_{1}|\psi_{2}^{(0)}\rangle = -\epsilon\langle y_{-}|\sigma_{x}|y_{-}\rangle = i\epsilon\langle y_{-}|y_{+}\rangle = 0. \\
\end{align}

Correções de segunda ordem para a energia
(\(E_{n}^{(2)} = \sum_{m\ne n}\frac{|\langle\psi_{n}^{(0)}|H_{1}|\psi_{m}^{(0)}\rangle|^{2}}{E_{n}^{(0)}-E_{m}^{(0)}}\)):

\begin{align}
E_{1}^{(2)} & = \sum_{m\ne 1}\frac{|\langle\psi_{1}^{(0)}|H_{1}|\psi_{m}^{(0)}\rangle|^{2}}{E_{1}^{(0)}-E_{m}^{(0)}} = \frac{|\langle\psi_{1}^{(0)}|H_{1}|\psi_{2}^{(0)}\rangle|^{2}}{E_{1}^{(0)}-E_{2}^{(0)}} = \frac{\epsilon^{2}|\langle y_{+}|\sigma_{x}|y_{-}\rangle|^{2}}{-\omega-\omega} = \frac{-\epsilon^{2}}{2\omega}, \\
E_{2}^{(2)} & =  \sum_{m\ne 2}\frac{|\langle\psi_{2}^{(0)}|H_{1}|\psi_{m}^{(0)}\rangle|^{2}}{E_{2}^{(0)}-E_{m}^{(0)}} = \frac{|\langle\psi_{2}^{(0)}|H_{1}|\psi_{1}^{(0)}\rangle|^{2}}{E_{2}^{(0)}-E_{1}^{(0)}} = \frac{\epsilon^{2}|\langle y_{-}|\sigma_{x}|y_{+}\rangle|^{2}}{\omega-(-\omega)} = \frac{\epsilon^{2}}{2\omega}.
\end{align}

Correções de primeira ordem para o estado
(\(|\psi_{n}^{(1)}\rangle = \sum_{m\ne n}\frac{\langle\psi_{m}^{(0)}|H_{1}|\psi_{n}^{(0)}\rangle}{E_{n}^{(0)}-E_{m}^{(0)}}|\psi_{m}^{(0)}\rangle\)):

\begin{align}
|\psi_{1}^{(1)}\rangle & = \sum_{m\ne 1}\frac{\langle\psi_{m}^{(0)}|H_{1}|\psi_{1}^{(0)}\rangle}{E_{1}^{(0)}-E_{m}^{(0)}}|\psi_{m}^{(0)}\rangle = \frac{\langle\psi_{2}^{(0)}|H_{1}|\psi_{1}^{(0)}\rangle}{E_{1}^{(0)}-E_{2}^{(0)}}|\psi_{2}^{(0)}\rangle = \frac{-\epsilon\langle y_{-}|\sigma_{x}|y_{+}\rangle}{-\omega-\omega}|y_{-}\rangle = \frac{i\epsilon}{2\omega}|y_{-}\rangle, \\
|\psi_{2}^{(1)}\rangle & = \sum_{m\ne 2}\frac{\langle\psi_{m}^{(0)}|H_{1}|\psi_{2}^{(0)}\rangle}{E_{2}^{(0)}-E_{m}^{(0)}}|\psi_{m}^{(0)}\rangle = \frac{\langle\psi_{1}^{(0)}|H_{1}|\psi_{2}^{(0)}\rangle}{E_{2}^{(0)}-E_{1}^{(0)}}|\psi_{1}^{(0)}\rangle = \frac{-\epsilon\langle y_{+}|\sigma_{x}|y_{-}\rangle}{\omega+\omega}|y_{+}\rangle = \frac{i\epsilon}{2\omega}|y_{+}\rangle.
\end{align}


    % Add a bibliography block to the postdoc
    
    
    
    \end{document}
